%======================================================================
\section{Track FND --- Foundation and upstream}\label{sec:fnd}
%======================================================================

Every other track depends on a fork that builds reproducibly, ships, and keeps up with FreeCAD. Two
foundation decisions changed in the interview. BruceWorx now tracks \textbf{upstream main} instead of
waiting for the 26.3 tag. Fixes found while reviewing the UI now have an epic of their own.

\begin{epic}{FND-01}{Build \& CI on Linux and macOS arm64}{M0}{\stProg}{FCB-P0-CI-01..06, P0-01, P0-06..08, P0-10}
Pixi environments and lockfiles for both platforms, ccache in CI, and the self-hosted Apple-silicon runner.
\file{BUILDING-BRUCEWORX.md} and \file{.github/workflows/bruceworx-ci.yml} exist, and the Linux STEP check
and hosted macOS lock are merged (\#4).
\emph{Acceptance:} a fresh clone builds on both platforms using only the documented commands, and CI is green
on both.
\end{epic}

\begin{epic}{FND-02}{Packaging, branding \& releases}{M0}{\stProg}{FCB-P0-02, P0-03, P0-09, P0-11, P0-17, P0-18}
The version string already reads \code{1.1.4-bw.1}, and the unsigned arm64 \code{.dmg} script is merged
(\#3). Remaining: tag a release carrying the dmg plus a source tarball, and decide on signing and
notarisation before anyone outside the team installs it.
\emph{Acceptance:} the tagged release installs on a clean Mac and opens a Fusion STEP export
(\code{LE-01001}, 200$\times$200$\times$8 envelope).
\end{epic}

\begin{epic}{FND-03}{Track upstream FreeCAD main}{M0}{\stTodo}{FCB-P1-16..21, P0-12 (replaced)}
\file{UPSTREAM-SYNC.md} currently says ``stay on 1.1.4, rebase onto 26.3 when it ships''. That changes to
continuous tracking:
\begin{itemize}
  \item \code{upstream-main} stays a fast-forward mirror. A nightly job rebases the \code{[BW]} series onto it
        and runs the BruceWorx tests plus the headless acceptance corpus (\eref{FND-05}).
  \item A conflict or red build opens an issue automatically. The \emph{conflict budget} is the number of
        upstream files our series touches. It is reported each night and must not grow without a reason
        recorded in the PR.
  \item Product releases are still cut from a tested point (\code{main} after a green nightly), not from
        arbitrary upstream commits.
\end{itemize}
\emph{Why:} the Assembly, TechDraw and FEM work that parity depends on lands upstream continuously.
Rebasing once a year makes every conflict expensive at once.
\emph{Acceptance:} the nightly rebase runs for two weeks; each conflict is resolved within a working day.
\end{epic}

\begin{epic}{FND-04}{Licences \& notices}{M0}{\stTodo}{FCB-P0-13}
FreeCAD is LGPL-2.1-or-later. The roadmap bundles more components: SheetMetal (LGPL-2.1-or-later), OCP, CadQuery and build123d
(Apache-2.0), possibly an OpenSCAD binary (GPL-2, shipped as a separate executable) and a fastener library
(licence to check \verify). Each needs a line in \file{NOTICE} and a third-party licence file in the app
bundle.
\emph{Acceptance:} every bundled component is listed; each release has a source tarball attached.
\end{epic}

\begin{epic}{FND-05}{Acceptance corpus \& test harness}{M0}{\stTodo}{FCB-P1-09, P3-08}
One corpus serves every epic: native and STEP versions of LE parts and assemblies, the 2-DOF arm, the quad
frame, and the benchmark models in this document. A single command runs every headless check.
Screenshots come from the existing before/after script. The renders in this document already come from
headless FreeCAD (\file{scripts/}), and the harness generalises that.
\emph{Acceptance:} \code{pixi run acceptance} runs all checks in CI in under 15 minutes.
\end{epic}

\begin{epic}{FND-06}{Fixes found in review}{M0}{\stProg}{(new)}
\begin{enumerate}
  \item \textbf{Linear Pattern direction 2 drops a row} (\S\ref{sec:bug}). \code{Spacings2} is declared
        \code{(\{\})}, which stores \code{[0.0]}. Fix the declaration, migrate saved files, add a test.
  \item \textbf{Workbench trim not applied.} The pack lists FEM, CAM, BIM, Robot and others as disabled
        (\file{bruceworx\_pack.py:118}), but the screenshot shows 13 tabs. The uncommitted
        \code{param\_fallback} change suggests the pack crashed on \code{int(None)} \verify.
  \item \textbf{Overlay model tree missing} on the screenshot, although it is on by default; root cause unknown \verify.
\end{enumerate}
\emph{Acceptance:} a $4\times3$ pattern gives 12 instances; only intended workbenches are visible; the tree shows on first run.
\end{epic}

\begin{epic}{FND-07}{On by default, each with an off switch}{M0}{\stDone}{FCB-P1-01, P1-10}
Merged in \#9. It remains a rule for every epic below: a feature ships enabled and can be switched off on
Preferences~$\rightarrow$~BruceWorx, and switching it off restores stock behaviour.
\end{epic}
